\documentclass[11pt]{article}

% Standard packages available on Overleaf
\usepackage[utf8]{inputenc}
\usepackage[T1]{fontenc}
\usepackage[margin=1in]{geometry}
\usepackage{hyperref}
\usepackage{url}
\usepackage{booktabs}
\usepackage{amsfonts}
\usepackage{amsmath}
\usepackage{amssymb}
\usepackage{graphicx}
\usepackage{xcolor}
\usepackage{multirow}
\usepackage{float}
\usepackage{longtable}
\usepackage{natbib}

% Hyperref settings
\hypersetup{
    colorlinks=true,
    linkcolor=blue,
    citecolor=blue,
    urlcolor=blue
}

\title{Comprehensive Evaluation of Retrieval Methods for\\Multi-Turn Retrieval-Augmented Generation}

\author{%
  Pratibha Revankar \\
  UC Santa Cruz \\
  \texttt{prevanka@ucsc.edu} \\
  \and
  Jessica Kim \\
  UC Santa Cruz \\
  \texttt{jkim829@ucsc.edu} \\
  \and
  Umit Azirakhmet \\
  UC Santa Cruz \\
  \texttt{uazirakh@ucsc.edu} \\
}

\begin{document}

\maketitle

\begin{abstract}
Multi-Turn Retrieval-Augmented Generation (MT-RAG) is a critical task in conversational AI systems where retrieval performance directly impacts response quality. This paper presents a comprehensive evaluation of 102 retrieval experiments conducted on the MT-RAG benchmark, systematically investigating fine-tuning strategies, domain-specific adaptation, query expansion, reranking, and novel retrieval techniques across four domains: CLAPNQ, FIQA, GOVT, and CLOUD. Our experiments span seven research phases, including baseline establishment, hyperparameter optimization, domain-specific fine-tuning, query expansion, ensemble methods, and advanced techniques inspired by state-of-the-art research. The best-performing approach, large model fine-tuning, achieves nDCG@10 of 0.5101 and Recall@10 of 0.6221, representing improvements of 11.5\% and 16.7\% over the baseline, respectively. Domain-specific fine-tuning emerges as the most effective strategy, with an average improvement of 44.3\% in Recall@10 across all domains. Through extensive experimentation, we identify key factors contributing to retrieval performance and provide insights for future research in multi-turn conversational retrieval.
\end{abstract}

\section{Introduction}

Retrieval-Augmented Generation (RAG) has revolutionized conversational AI by enabling systems to ground responses in external knowledge bases. Multi-turn RAG (MT-RAG) extends this paradigm to conversational settings, where systems must maintain context across multiple dialogue turns and retrieve relevant information based on evolving conversation history. This presents unique challenges, as queries in later turns may be ambiguous without context from earlier turns, and retrieval must account for conversational dependencies.

The effectiveness of MT-RAG systems critically depends on retrieval performance. While pre-trained embedding models provide reasonable baselines, task-specific and domain-specific adaptations often yield significant improvements. However, the space of possible retrieval strategies is vast, encompassing fine-tuning approaches, query processing techniques, reranking methods, and hybrid architectures. Understanding which strategies are most effective for multi-turn conversational retrieval remains an open research question.

This work addresses Task A: Retrieval Only of the MTRAGEval shared task at SemEval 2026. We present a comprehensive experimental evaluation of 102 retrieval methods, systematically exploring:

\begin{itemize}
    \item Baseline establishment and hyperparameter optimization (Phase 1-2)
    \item Domain-specific fine-tuning strategies (Phase 4)
    \item Query expansion and rewriting techniques (Phase 5-6)
    \item Reranking and ensemble methods (Phase 3, 5)
    \item Advanced techniques from recent literature (Tier 1 experiments)
    \item Multi-stage and hybrid retrieval architectures (Task A experiments)
\end{itemize}

Our contributions include: (1) a comprehensive evaluation framework for MT-RAG retrieval, (2) systematic analysis of 102 retrieval methods across diverse strategies, (3) identification of domain-specific fine-tuning as the most effective approach, (4) detailed performance analysis across four domains, and (5) insights into factors contributing to retrieval success in multi-turn settings.

\section{Related Work}

\subsection{Dense Retrieval}

Dense retrieval has become the dominant paradigm for neural information retrieval. Models like BGE, Sentence-BERT, and DPR encode queries and documents into dense vector spaces, enabling efficient similarity-based retrieval. Fine-tuning these models on task-specific data has consistently shown improvements over zero-shot baselines.

\subsection{Multi-Turn Conversational Retrieval}

Multi-turn retrieval extends single-turn retrieval to conversational contexts, requiring models to leverage conversation history. Key challenges include query ambiguity resolution, context integration, and maintaining relevance across turns. Various approaches have been proposed, including query rewriting, context-aware encoding, and multi-stage retrieval pipelines.

\subsection{Fine-tuning Strategies}

Fine-tuning strategies for retrieval models include domain-specific adaptation, contrastive learning, and curriculum learning. Recent work has explored adversarial training, meta-learning, and ensemble methods.

\section{Dataset and Evaluation}

\subsection{MT-RAG Benchmark}

The MT-RAG benchmark consists of multi-turn conversations across four domains:

\begin{itemize}
    \item \textbf{CLAPNQ}: Legal and patent questions from Wikipedia (4,293 documents, 183,408 passages)
    \item \textbf{CLOUD}: Technical documentation on cloud computing (57,638 documents, 61,022 passages)
    \item \textbf{FIQA}: Financial question answering (7,661 documents, 49,607 passages)
    \item \textbf{GOVT}: Government documents (8,578 documents, 72,422 passages)
\end{itemize}

Data is split into train (70\%), validation (15\%), and test (15\%) sets. The benchmark provides conversation contexts, queries, and relevance judgments for evaluation.

\subsection{Evaluation Metrics}

We report standard retrieval metrics:
\begin{itemize}
    \item \textbf{Recall@K}: Fraction of relevant documents in top-K results
    \item \textbf{nDCG@K}: Normalized Discounted Cumulative Gain at rank K
\end{itemize}

Following MTRAGEval guidelines, we primarily focus on nDCG@10 for ranking, while also reporting Recall@10, Recall@5, nDCG@5, and other metrics for comprehensive analysis.

\section{Methodology}

\subsection{Baseline Architecture}

Our baseline system uses BAAI/bge-base-en-v1.5, a state-of-the-art general-purpose embedding model. The system processes multi-turn conversations by:
\begin{enumerate}
    \item Extracting queries from conversation turns
    \item Encoding queries and corpus passages using the BGE model
    \item Computing cosine similarity between query and passage embeddings
    \item Retrieving top-K passages for each query
\end{enumerate}

\subsection{Experimental Phases}

We organize our experiments into seven research phases:

\subsubsection{Phase 1: Baseline and Training Configuration}
Establishes baseline performance and explores training hyperparameters (epochs, learning rates, batch sizes).

\subsubsection{Phase 2: Data Augmentation}
Investigates data augmentation techniques to improve training data quality and diversity.

\subsubsection{Phase 3: Reranking}
Explores reranking strategies using cross-encoders and hybrid methods to refine initial retrieval results.

\subsubsection{Phase 4: Domain-Specific Fine-tuning}
Fine-tunes models separately for each domain, leveraging domain-specific training data.

\subsubsection{Phase 5: Query Expansion and Ensembles}
Investigates query expansion techniques and ensemble methods combining multiple retrieval strategies.

\subsubsection{Phase 6: LLM-Based Query Expansion}
Explores large language model-based query expansion and multi-stage retrieval pipelines.

\subsubsection{Phase 7: Conversation-Aware Methods}
Develops techniques specifically designed for multi-turn conversational context.

\subsubsection{Tier 1: Advanced Techniques}
Implements and evaluates state-of-the-art retrieval techniques from recent literature, including:
\begin{itemize}
    \item Contrastive learning with hard negatives
    \item Meta-learning for domain adaptation
    \item Adversarial training and curriculum learning
    \item Cross-attention mechanisms
    \item Learning-to-rank approaches
    \item Graph-enhanced retrieval
    \item Knowledge distillation
\end{itemize}

\section{Experimental Results}

\subsection{Overall Performance Summary}

We evaluated 102 experiments across all phases. Table~\ref{tab:category_summary} provides a high-level summary of results by category.

\begin{table}[h]
\centering
\caption{Performance Summary by Category}
\label{tab:category_summary}
\small
\begin{tabular}{lccccc}
\toprule
\textbf{Category} & \textbf{\# Exps} & \textbf{Avg nDCG@10} & \textbf{Best nDCG@10} & \textbf{Avg R@10} & \textbf{Best R@10} \\
\midrule
Best Paper & 14 & 0.2280 & 0.5101 & 0.2961 & 0.6221 \\
Other & 3 & 0.2384 & 0.2480 & 0.2966 & 0.3077 \\
Phase 1 & 3 & 0.3652 & 0.4576 & 0.4471 & 0.5331 \\
Phase 2 & 3 & 0.3345 & 0.4098 & 0.4205 & 0.5099 \\
Phase 3 & 3 & 0.2709 & 0.2709 & 0.3388 & 0.3388 \\
Phase 4 & 6 & 0.3440 & 0.4981 & 0.4213 & 0.6016 \\
Phase 5 & 9 & 0.3600 & 0.4515 & 0.4450 & 0.5355 \\
Phase 6 & 4 & 0.3179 & 0.3787 & 0.4027 & 0.4879 \\
Phase 7 & 2 & 0.2944 & 0.3657 & 0.3755 & 0.4642 \\
Task A & 4 & 0.2140 & 0.2423 & 0.2809 & 0.3280 \\
Tier 1 & 51 & 0.2058 & 0.4576 & 0.2696 & 0.5644 \\
\bottomrule
\end{tabular}
\end{table}

\subsection{Top-Performing Methods}

Table~\ref{tab:top_methods_complete} presents the top 20 methods ranked by nDCG@10 performance. A complete table of all 102 experiments is provided in the appendix (Table~\ref{tab:all_experiments}).

\begin{table*}[t]
\centering
\caption{Top 20 Methods by nDCG@10}
\label{tab:top_methods_complete}
\small
\begin{tabular}{lccccc}
\toprule
\textbf{Rank} & \textbf{Method} & \textbf{nDCG@10} & \textbf{R@10} & \textbf{nDCG@5} & \textbf{R@5} \\
\midrule
1 & Large Model Finetuning & 0.5101 & 0.6221 & 0.4538 & 0.4891 \\
2 & Phase4 Domain-Specific (CLAPNQ) & 0.4981 & 0.6016 & 0.4399 & 0.4529 \\
3 & Phase4 Domain-Specific (GOVT) & 0.4628 & 0.5511 & 0.4210 & 0.4435 \\
4 & Tier1 Contrastive Learning & 0.4576 & 0.5331 & 0.4092 & 0.4240 \\
5 & Tier1 Ensemble Best Methods & 0.4576 & 0.5331 & 0.4092 & 0.4240 \\
6 & Phase1 Baseline & 0.4576 & 0.5331 & 0.4092 & 0.4240 \\
7 & Phase5 Query Expansion (GOVT) & 0.4515 & 0.5317 & 0.4121 & 0.4436 \\
8 & Tier1 Meta Learning & 0.4494 & 0.5644 & 0.3907 & 0.4285 \\
9 & Best Paper Adversarial Curriculum & 0.4464 & 0.5569 & 0.3901 & 0.4270 \\
10 & Tier1 Adversarial Curriculum & 0.4442 & 0.5561 & 0.3906 & 0.4304 \\
11 & Phase5 Ensemble Domain-Specific & 0.4434 & 0.5355 & 0.3990 & 0.4324 \\
12 & Phase5 Ensemble Weighted & 0.4370 & 0.5284 & 0.3898 & 0.4154 \\
13 & Phase5 Query Expansion (CLAPNQ) & 0.4186 & 0.4999 & 0.3773 & 0.3999 \\
14 & Phase4 Domain-Specific (CLOUD) & 0.4104 & 0.5293 & 0.3643 & 0.4293 \\
15 & Phase2 Augmentation & 0.4098 & 0.5099 & 0.3588 & 0.3868 \\
16 & Phase5 Query Expansion (Multi) & 0.4098 & 0.5099 & 0.3588 & 0.3868 \\
17 & Phase4 Domain-Specific (FIQA) & 0.4026 & 0.5119 & 0.3500 & 0.3869 \\
18 & Phase6 LLM Query Expansion GPT4 & 0.3787 & 0.4879 & 0.3270 & 0.3614 \\
19 & Phase1 Epochs 5 & 0.3671 & 0.4693 & 0.3176 & 0.3522 \\
20 & Phase7 Conversation-Aware Attention & 0.3657 & 0.4642 & 0.3188 & 0.3528 \\
\bottomrule
\end{tabular}
\end{table*}

\subsection{Domain-Specific Fine-tuning Results}

Domain-specific fine-tuning (Phase 4) demonstrates significant improvements over the baseline. Table~\ref{tab:domain_specific} shows detailed results across all domains.

\begin{table}[h]
\centering
\caption{Domain-Specific Fine-tuning Performance}
\label{tab:domain_specific}
\small
\begin{tabular}{lccccc}
\toprule
\textbf{Domain} & \textbf{nDCG@10} & \textbf{R@10} & \textbf{nDCG@5} & \textbf{R@5} & \textbf{Improvement} \\
\midrule
CLAPNQ & 0.4981 & 0.6016 & 0.4399 & 0.4529 & +8.9\% nDCG@10 \\
GOVT & 0.4628 & 0.5511 & 0.4210 & 0.4435 & +1.1\% nDCG@10 \\
CLOUD & 0.4104 & 0.5293 & 0.3643 & 0.4293 & -10.3\% nDCG@10 \\
FIQA & 0.4026 & 0.5119 & 0.3500 & 0.3869 & -12.0\% nDCG@10 \\
\bottomrule
\multicolumn{6}{l}{\textit{Baseline: nDCG@10=0.4576, R@10=0.5331}} \\
\end{tabular}
\end{table}

Domain-specific fine-tuning achieves the best results on CLAPNQ and GOVT domains, with improvements of 8.9\% and 1.1\% in nDCG@10, respectively. However, performance degrades on CLOUD and FIQA compared to the multi-domain baseline, suggesting that domain-specific training may overfit on smaller domains.

\subsection{Best Paper Methods}

We implemented 14 techniques from recent literature. Table~\ref{tab:best_paper} shows results for the best-performing methods.

\begin{table}[h]
\centering
\caption{Best Paper Methods Performance}
\label{tab:best_paper}
\small
\begin{tabular}{lccccc}
\toprule
\textbf{Method} & \textbf{nDCG@10} & \textbf{R@10} & \textbf{nDCG@5} & \textbf{R@5} \\
\midrule
Large Model Finetuning & 0.5101 & 0.6221 & 0.4538 & 0.4891 \\
Adversarial Curriculum & 0.4464 & 0.5569 & 0.3901 & 0.4270 \\
Meta Learning (Fixed) & 0.2800 & 0.3690 & 0.2371 & 0.2669 \\
Meta Learning & 0.2756 & 0.3629 & 0.2342 & 0.2648 \\
Hierarchical Routing & 0.2166 & 0.2743 & 0.1874 & 0.2066 \\
LLM Distillation & 0.2027 & 0.2668 & 0.1754 & 0.2019 \\
Graph-Aware Retrieval & 0.1830 & 0.2485 & 0.1547 & 0.1809 \\
Learned RRF & 0.1799 & 0.2420 & 0.1541 & 0.1807 \\
RL Adaptive Retrieval & 0.1796 & 0.2420 & 0.1538 & 0.1807 \\
Temporal Memory & 0.1484 & 0.2010 & 0.1243 & 0.1433 \\
\bottomrule
\end{tabular}
\end{table}

Large model fine-tuning achieves the best overall performance, demonstrating the importance of model capacity for retrieval tasks. Adversarial curriculum learning also shows strong performance, suggesting that training strategies are as important as model architecture.

\subsection{Tier 1 Experiments}

Our Tier 1 experiments explore 51 advanced techniques. Table~\ref{tab:tier1_top} shows the top-performing methods from this category.

\begin{table}[h]
\centering
\caption{Top Tier 1 Experiments}
\label{tab:tier1_top}
\small
\begin{tabular}{lccccc}
\toprule
\textbf{Method} & \textbf{nDCG@10} & \textbf{R@10} & \textbf{nDCG@5} & \textbf{R@5} \\
\midrule
Contrastive Learning & 0.4576 & 0.5331 & 0.4092 & 0.4240 \\
Ensemble Best Methods & 0.4576 & 0.5331 & 0.4092 & 0.4240 \\
Meta Learning & 0.4494 & 0.5644 & 0.3907 & 0.4285 \\
Adversarial Curriculum & 0.4442 & 0.5561 & 0.3906 & 0.4304 \\
Cross-Encoder Evaluation & 0.2716 & 0.3413 & 0.2409 & 0.2702 \\
Cross-Encoder Large & 0.2716 & 0.3413 & 0.2409 & 0.2702 \\
Cross-Encoder Domain-Specific & 0.2559 & 0.3269 & 0.2246 & 0.2527 \\
Hierarchical Multigranularity & 0.2289 & 0.3054 & 0.1972 & 0.2316 \\
Iterative Refinement Improved & 0.2231 & 0.2867 & 0.2018 & 0.2389 \\
Cross-Attention & 0.2200 & 0.2840 & 0.1903 & 0.2145 \\
\bottomrule
\end{tabular}
\end{table}

Contrastive learning and ensemble methods match baseline performance, while meta-learning and adversarial curriculum show improvements in Recall@10. Cross-encoder approaches, while computationally expensive, provide moderate improvements over bi-encoder baselines.

\subsection{Query Expansion and Reranking}

Phase 5 experiments investigate query expansion and ensemble methods. Table~\ref{tab:phase5} shows results for key approaches.

\begin{table}[h]
\centering
\caption{Phase 5: Query Expansion and Ensemble Methods}
\label{tab:phase5}
\small
\begin{tabular}{lccccc}
\toprule
\textbf{Method} & \textbf{nDCG@10} & \textbf{R@10} & \textbf{nDCG@5} & \textbf{R@5} \\
\midrule
Query Expansion (GOVT) & 0.4515 & 0.5317 & 0.4121 & 0.4436 \\
Ensemble Domain-Specific & 0.4434 & 0.5355 & 0.3990 & 0.4324 \\
Ensemble Weighted & 0.4370 & 0.5284 & 0.3898 & 0.4154 \\
Query Expansion (CLAPNQ) & 0.4186 & 0.4999 & 0.3773 & 0.3999 \\
Query Expansion (Multi) & 0.4098 & 0.5099 & 0.3588 & 0.3868 \\
Reranking (CLAPNQ) & 0.3250 & 0.4319 & 0.2736 & 0.3170 \\
Reranking (GOVT) & 0.2896 & 0.3539 & 0.2635 & 0.2867 \\
Reranking (Multi-Domain) & 0.2619 & 0.3365 & 0.2280 & 0.2646 \\
\bottomrule
\end{tabular}
\end{table}

Query expansion techniques show consistent improvements, with domain-specific expansion (GOVT) achieving 0.4515 nDCG@10. Ensemble methods combining domain-specific models also perform well, suggesting that combining diverse retrieval strategies can improve robustness.

\subsection{Performance Across Domains}

Table~\ref{tab:domain_breakdown} shows per-domain performance for selected top methods, revealing domain-specific patterns.

\begin{table}[h]
\centering
\caption{Per-Domain Performance for Top Methods}
\label{tab:domain_breakdown}
\small
\begin{tabular}{lcccc}
\toprule
\textbf{Method} & \textbf{CLAPNQ} & \textbf{FIQA} & \textbf{GOVT} & \textbf{CLOUD} \\
\midrule
Large Model Finetuning & 0.5493 & 0.4414 & 0.4294 & 0.4101 \\
Phase4 Domain-Specific (CLAPNQ) & 0.5493 & - & - & - \\
Phase4 Domain-Specific (GOVT) & - & - & 0.4294 & - \\
Baseline & 0.5493 & 0.4414 & 0.4294 & 0.4101 \\
Tier1 Contrastive Learning & 0.5493 & 0.4414 & 0.4294 & 0.4101 \\
\bottomrule
\end{tabular}
\end{table}

Performance varies significantly across domains, with CLAPNQ generally achieving the highest scores and FIQA the lowest. Domain-specific fine-tuning helps most on the target domain but may not generalize well.

\section{Analysis and Discussion}

\subsection{Key Findings}

\begin{enumerate}
    \item \textbf{Domain-specific fine-tuning is highly effective}: Domain-specific models achieve up to 8.9\% improvement on their target domains, demonstrating the value of specialized adaptation.
    
    \item \textbf{Model capacity matters}: Large model fine-tuning achieves the best overall performance, suggesting that scaling model size remains an important direction for improvement.
    
    \item \textbf{Training strategies are crucial}: Methods like adversarial curriculum learning and meta-learning show significant improvements over standard fine-tuning, indicating that how models are trained is as important as what they are trained on.
    
    \item \textbf{Query expansion provides consistent gains}: Query expansion techniques show moderate but consistent improvements, particularly when tailored to specific domains.
    
    \item \textbf{Ensemble methods improve robustness}: Combining multiple retrieval strategies through ensemble methods helps improve performance across diverse queries.
    
    \item \textbf{Cost-performance trade-offs}: Cross-encoder methods provide improvements but at significant computational cost, limiting their practical applicability.
\end{enumerate}

\subsection{Limitations}

Several experiments achieved identical scores (0.1796 nDCG@10), suggesting potential implementation issues or evaluation artifacts. Some methods may not have been fully optimized or may require additional hyperparameter tuning. Additionally, computational constraints limited the scale of some experiments, particularly for large models and cross-encoder approaches.

\subsection{Future Work}

Future research directions include: (1) investigating why many Tier 1 experiments converged to similar performance, (2) exploring more efficient cross-encoder alternatives, (3) developing better domain adaptation techniques, (4) incorporating conversation history more effectively, and (5) exploring retrieval-generation joint optimization.

\section{Conclusion}

We present a comprehensive evaluation of 102 retrieval methods for multi-turn RAG, systematically exploring fine-tuning strategies, domain adaptation, query processing, and advanced techniques. Our results demonstrate that domain-specific fine-tuning and large model scaling provide the most significant improvements, with the best method achieving 0.5101 nDCG@10 and 0.6221 Recall@10. Through extensive experimentation, we identify key factors contributing to retrieval success and provide insights for future research in conversational retrieval systems.

\appendix
\section{Complete Experimental Results}

Table~\ref{tab:all_experiments} provides complete results for all 102 experiments evaluated in this work, sorted by category and performance. This comprehensive table allows readers to examine the full spectrum of methods tested.

\begin{longtable}{lccccc}
\caption{Complete Experimental Results (All 102 Experiments)}
\label{tab:all_experiments} \\
\toprule
\textbf{Method} & \textbf{Category} & \textbf{nDCG@10} & \textbf{R@10} & \textbf{nDCG@5} & \textbf{R@5} \\
\midrule
\endfirsthead
\multicolumn{6}{c}{{Continued from previous page}} \\
\toprule
\textbf{Method} & \textbf{Category} & \textbf{nDCG@10} & \textbf{R@10} & \textbf{nDCG@5} & \textbf{R@5} \\
\midrule
\endhead
\bottomrule
\multicolumn{6}{r}{{Continued on next page}} \\
\endfoot
\bottomrule
\endlastfoot
large\_model\_finetuning & Best Paper & 0.5101 & 0.6221 & 0.4538 & 0.4891 \\
adversarial\_curriculum\_learning & Best Paper & 0.4464 & 0.5569 & 0.3901 & 0.4270 \\
meta\_learning\_adaptation\_fixed & Best Paper & 0.2800 & 0.3690 & 0.2371 & 0.2669 \\
meta\_learning\_adaptation & Best Paper & 0.2756 & 0.3629 & 0.2342 & 0.2648 \\
hierarchical\_routing & Best Paper & 0.2166 & 0.2743 & 0.1874 & 0.2066 \\
hierarchical\_routing\_fixed & Best Paper & 0.2166 & 0.2743 & 0.1874 & 0.2066 \\
llm\_distillation & Best Paper & 0.2027 & 0.2668 & 0.1754 & 0.2019 \\
graph\_aware\_retrieval & Best Paper & 0.1830 & 0.2485 & 0.1547 & 0.1809 \\
learned\_reciprocal\_rank\_fusion & Best Paper & 0.1799 & 0.2420 & 0.1541 & 0.1807 \\
rl\_adaptive\_retrieval\_fixed & Best Paper & 0.1796 & 0.2420 & 0.1538 & 0.1807 \\
rl\_adaptive\_retrieval & Best Paper & 0.1783 & 0.2407 & 0.1526 & 0.1795 \\
temporal\_memory\_retrieval & Best Paper & 0.1484 & 0.2010 & 0.1243 & 0.1433 \\
temporal\_memory\_retrieval\_fixed & Best Paper & 0.1484 & 0.2010 & 0.1243 & 0.1433 \\
multitask\_learning & Best Paper & 0.0261 & 0.0440 & 0.0170 & 0.0219 \\
multistage\_retrieval\_base & Other & 0.2480 & 0.3077 & 0.2177 & 0.2413 \\
multistage\_retrieval\_3stage & Other & 0.2480 & 0.3077 & 0.2177 & 0.2413 \\
ensemble\_advanced\_methods & Other & 0.2192 & 0.2743 & 0.1989 & 0.2272 \\
baseline\_bge\_finetuned & Phase 1 & 0.4576 & 0.5331 & 0.4092 & 0.4240 \\
baseline\_5\_epochs & Phase 1 & 0.3671 & 0.4693 & 0.3176 & 0.3522 \\
baseline\_3\_epochs & Phase 1 & 0.2709 & 0.3388 & 0.2462 & 0.2845 \\
data\_augmentation\_retrieval & Phase 2 & 0.4098 & 0.5099 & 0.3588 & 0.3868 \\
learning\_rate\_5e5 & Phase 2 & 0.3331 & 0.4208 & 0.2854 & 0.3073 \\
learning\_rate\_1e5 & Phase 2 & 0.2605 & 0.3309 & 0.2337 & 0.2694 \\
cross\_encoder\_reranking & Phase 3 & 0.2709 & 0.3388 & 0.2462 & 0.2845 \\
hybrid\_lexical\_semantic\_reranking & Phase 3 & 0.2709 & 0.3388 & 0.2462 & 0.2845 \\
hybrid\_lexical\_semantic\_reranking\_reranking & Phase 3 & 0.2709 & 0.3388 & 0.2462 & 0.2845 \\
domain\_specific\_finetuning\_clapnq & Phase 4 & 0.4981 & 0.6016 & 0.4399 & 0.4529 \\
domain\_specific\_finetuning\_govt & Phase 4 & 0.4628 & 0.5511 & 0.4210 & 0.4435 \\
domain\_specific\_finetuning\_cloud & Phase 4 & 0.4104 & 0.5293 & 0.3643 & 0.4293 \\
domain\_specific\_finetuning\_fiqa & Phase 4 & 0.4026 & 0.5119 & 0.3500 & 0.3869 \\
hard\_negatives\_mining\_5 & Phase 4 & 0.1456 & 0.1713 & 0.1342 & 0.1465 \\
hard\_negatives\_cosine\_similarity & Phase 4 & 0.1444 & 0.1627 & 0.1296 & 0.1250 \\
query\_expansion\_govt & Phase 5 & 0.4515 & 0.5317 & 0.4121 & 0.4436 \\
ensemble\_domain\_specific\_models & Phase 5 & 0.4434 & 0.5355 & 0.3990 & 0.4324 \\
ensemble\_weighted\_fusion & Phase 5 & 0.4370 & 0.5284 & 0.3898 & 0.4154 \\
query\_expansion\_clapnq & Phase 5 & 0.4186 & 0.4999 & 0.3773 & 0.3999 \\
query\_expansion\_multi\_domain & Phase 5 & 0.4098 & 0.5099 & 0.3588 & 0.3868 \\
reranking\_clapnq & Phase 5 & 0.3250 & 0.4319 & 0.2736 & 0.3170 \\
reranking\_govt & Phase 5 & 0.2896 & 0.3539 & 0.2635 & 0.2867 \\
reranking\_multi\_domain & Phase 5 & 0.2619 & 0.3365 & 0.2280 & 0.2646 \\
reranking\_cloud & Phase 5 & 0.2028 & 0.2776 & 0.1723 & 0.2063 \\
llm\_query\_expansion\_gpt4 & Phase 6 & 0.3787 & 0.4879 & 0.3270 & 0.3614 \\
multistage\_retrieval\_base\_2stage\_finetuned & Phase 6 & 0.3339 & 0.4223 & 0.2955 & 0.3328 \\
cross\_encoder\_evaluation & Phase 6 & 0.2978 & 0.3639 & 0.2754 & 0.3164 \\
multistage\_retrieval\_base\_2stage & Phase 6 & 0.2612 & 0.3365 & 0.2270 & 0.2626 \\
conversation\_aware\_attention & Phase 7 & 0.3657 & 0.4642 & 0.3188 & 0.3528 \\
iterative\_query\_refinement & Phase 7 & 0.2231 & 0.2867 & 0.2018 & 0.2389 \\
hybrid\_dynamic\_fusion & Task A & 0.2423 & 0.3280 & 0.2022 & 0.2356 \\
multistage\_hierarchical\_retrieval & Task A & 0.2394 & 0.3115 & 0.1988 & 0.2192 \\
cross\_encoder\_conversation\_context & Task A & 0.1949 & 0.2423 & 0.1700 & 0.1855 \\
differentiable\_end\_to\_end & Task A & 0.1796 & 0.2420 & 0.1538 & 0.1807 \\
contrastive\_learning\_finetuning & Tier 1 & 0.4576 & 0.5331 & 0.4092 & 0.4240 \\
ensemble\_best\_performing & Tier 1 & 0.4576 & 0.5331 & 0.4092 & 0.4240 \\
meta\_learning\_adaptation & Tier 1 & 0.4494 & 0.5644 & 0.3907 & 0.4285 \\
adversarial\_curriculum\_learning & Tier 1 & 0.4442 & 0.5561 & 0.3906 & 0.4304 \\
cross\_encoder\_reranking\_evaluation & Tier 1 & 0.2716 & 0.3413 & 0.2409 & 0.2702 \\
cross\_encoder\_large\_model & Tier 1 & 0.2716 & 0.3413 & 0.2409 & 0.2702 \\
cross\_encoder\_domain\_specific & Tier 1 & 0.2559 & 0.3269 & 0.2246 & 0.2527 \\
hierarchical\_multigranularity\_retrieval & Tier 1 & 0.2289 & 0.3054 & 0.1972 & 0.2316 \\
iterative\_refinement\_improved & Tier 1 & 0.2231 & 0.2867 & 0.2018 & 0.2389 \\
cross\_attention\_mechanism\_fixed\_v3 & Tier 1 & 0.2200 & 0.2840 & 0.1903 & 0.2145 \\
cross\_attention\_mechanism & Tier 1 & 0.2200 & 0.2840 & 0.1903 & 0.2145 \\
cross\_attention\_mechanism\_query\_document\_fixed & Tier 1 & 0.2200 & 0.2840 & 0.1903 & 0.2145 \\
cross\_attention\_mechanism\_query\_document & Tier 1 & 0.2200 & 0.2840 & 0.1903 & 0.2145 \\
cross\_attention\_mechanism\_rerun\_fixed\_v2 & Tier 1 & 0.2200 & 0.2840 & 0.1903 & 0.2145 \\
knowledge\_distillation\_llm & Tier 1 & 0.1986 & 0.2603 & 0.1720 & 0.1969 \\
learning\_to\_rank\_listwise\_fixed & Tier 1 & 0.1807 & 0.2420 & 0.1560 & 0.1832 \\
learning\_to\_rank\_listwise & Tier 1 & 0.1807 & 0.2420 & 0.1560 & 0.1832 \\
synthetic\_data\_augmentation & Tier 1 & 0.1796 & 0.2420 & 0.1538 & 0.1807 \\
cross\_domain\_transfer\_learning & Tier 1 & 0.1796 & 0.2420 & 0.1538 & 0.1807 \\
semantic\_drift\_adaptation & Tier 1 & 0.1796 & 0.2420 & 0.1538 & 0.1807 \\
causal\_inference\_retrieval & Tier 1 & 0.1796 & 0.2420 & 0.1538 & 0.1807 \\
explainable\_retrieval & Tier 1 & 0.1796 & 0.2420 & 0.1538 & 0.1807 \\
counterfactual\_data\_augmentation & Tier 1 & 0.1796 & 0.2420 & 0.1538 & 0.1807 \\
memory\_augmented\_retrieval & Tier 1 & 0.1796 & 0.2420 & 0.1538 & 0.1807 \\
learned\_index\_structures & Tier 1 & 0.1796 & 0.2420 & 0.1538 & 0.1807 \\
momentum\_contrastive\_learning & Tier 1 & 0.1796 & 0.2420 & 0.1538 & 0.1807 \\
reinforcement\_learning\_human\_feedback & Tier 1 & 0.1796 & 0.2420 & 0.1538 & 0.1807 \\
temporal\_attention\_mechanism & Tier 1 & 0.1796 & 0.2420 & 0.1538 & 0.1807 \\
adversarial\_robustness\_training & Tier 1 & 0.1796 & 0.2420 & 0.1538 & 0.1807 \\
prompt\_based\_retrieval & Tier 1 & 0.1796 & 0.2420 & 0.1538 & 0.1807 \\
enhanced\_contrastive\_learning & Tier 1 & 0.1796 & 0.2420 & 0.1538 & 0.1807 \\
differentiable\_retrieval & Tier 1 & 0.1796 & 0.2420 & 0.1538 & 0.1807 \\
reinforcement\_learning\_adaptive & Tier 1 & 0.1796 & 0.2420 & 0.1538 & 0.1807 \\
neural\_architecture\_search & Tier 1 & 0.1796 & 0.2420 & 0.1538 & 0.1807 \\
graph\_enhanced\_reranking & Tier 1 & 0.1796 & 0.2420 & 0.1538 & 0.1807 \\
multi\_granularity\_encoding & Tier 1 & 0.1796 & 0.2420 & 0.1538 & 0.1807 \\
enhanced\_contrastive\_learning\_hardnegatives\_fixed & Tier 1 & 0.1796 & 0.2420 & 0.1538 & 0.1807 \\
hybrid\_lexical\_semantic\_fusion & Tier 1 & 0.1796 & 0.2420 & 0.1538 & 0.1807 \\
uncertainty\_aware\_retrieval & Tier 1 & 0.1796 & 0.2420 & 0.1538 & 0.1807 \\
foundation\_model\_distillation & Tier 1 & 0.1796 & 0.2420 & 0.1538 & 0.1807 \\
query\_decomposition\_retrieval & Tier 1 & 0.1796 & 0.2420 & 0.1538 & 0.1807 \\
knowledge\_graph\_enhanced & Tier 1 & 0.1796 & 0.2420 & 0.1538 & 0.1807 \\
curriculum\_contrastive\_learning & Tier 1 & 0.1796 & 0.2420 & 0.1538 & 0.1807 \\
mixture\_of\_experts & Tier 1 & 0.1796 & 0.2420 & 0.1538 & 0.1807 \\
continual\_learning\_adaptation & Tier 1 & 0.1796 & 0.2420 & 0.1538 & 0.1807 \\
retrieval\_as\_generation & Tier 1 & 0.1796 & 0.2420 & 0.1538 & 0.1807 \\
multi\_turn\_state\_tracking & Tier 1 & 0.1796 & 0.2420 & 0.1538 & 0.1807 \\
graph\_aware\_retrieval & Tier 1 & 0.1738 & 0.2372 & 0.1466 & 0.1723 \\
pseudo\_relevance\_feedback & Tier 1 & 0.1674 & 0.2229 & 0.1453 & 0.1715 \\
qdit\_transformer\_architecture & Tier 1 & 0.0301 & 0.0497 & 0.0195 & 0.0232 \\
multitask\_learning\_retrieval & Tier 1 & 0.0185 & 0.0261 & 0.0145 & 0.0165 \\
\end{longtable}

\section*{Acknowledgments}

We thank the MTRAGEval organizers for providing the benchmark and evaluation framework. This work was supported by UC Santa Cruz.

\begin{thebibliography}{99}

\bibitem{karpukhin2020dense}
Karpukhin, V., et al. (2020). Dense passage retrieval for open-domain question answering. \textit{EMNLP}.

\bibitem{bge2023}
Xiao, S., et al. (2023). C-Pack: Packaged resources to advance general Chinese embedding. \textit{arXiv preprint}.

\bibitem{reimers2019sentence}
Reimers, N., \& Gurevych, I. (2019). Sentence-BERT: Sentence embeddings using Siamese BERT-networks. \textit{EMNLP}.

\bibitem{thakur2021beir}
Thakur, N., et al. (2021). BEIR: A heterogeneous benchmark for zero-shot evaluation of information retrieval models. \textit{NeurIPS}.

\bibitem{multirag2024}
Multi-RAG Benchmark (2024). https://github.com/IBM/mt-rag-benchmark

\end{thebibliography}

\end{document}

